\documentclass[10pt,a4paper]{amsart}
\usepackage[margin=19mm]{geometry}
\usepackage{amsmath,amssymb,mathtools}
\usepackage[scr=rsfs,cal=euler]{mathalfa}
\usepackage{xfrac}
\usepackage[unicode,hidelinks,bookmarks=false]{hyperref}
\newcommand{\ie}{i.e.\ }
\newcommand{\cf}{cf.\ }
\newcommand{\resp}{respectively\ }
\newcommand{\Subsection}{\subsection}
\providecommand{\nobreakdash}{\nobreak\hbox{-}}
\setlength{\emergencystretch}{2em}
\multlinegap0pt
\usepackage{tikz}
\usepackage{amssymb}
\usepackage{mathtools}
\newtheorem{lem}{Lemma}
\title{Middle convolution for Lie algebra representations}
\author{Kazuki Hiroe}
\date{Source: arXiv:2605.09828v2 (CC BY 4.0)}
\begin{document}
\maketitle

\noindent\emph{Source and licence.} Middle convolution for Lie algebra representations. Version arXiv:2605.09828v2 (CC BY 4.0), distributed by arXiv under the Creative Commons Attribution 4.0 International licence, http://creativecommons.org/licenses/by/4.0/, as recorded in the arXiv licence metadata for this exact version and shown on the abstract page https://arxiv.org/abs/2605.09828v2. Full licence notice: \texttt{licenses/arxiv-2605.09828-CC-BY-4.0.txt}.\par\smallskip

\noindent\textit{Editorial scope.} Counted unit: the article's lem lem (source0.tex lines 863--870). The article's complete original proof of it is delivered with it (source0.tex lines 871--935, one \texttt{proof} environment of 65 lines). No mathematics is written, repaired or re-derived: the statement and the proof are the article's own lines, in the article's own notation, and no retained line is substituted or re-flowed.  No further source-local context is needed: every cross-reference of every retained line resolves inside the two delivered spans.  Nothing was omitted from any retained span, so the package carries no omission record.  No retained line contains a citation, so the wrapper carries no bibliography and no bibliography entry.  No retained line contains a figure environment, an image-inclusion command or drawing code, so no image is delivered and the package declares no asset dependency.  Editorial formatting only, in addition: an AMS \texttt{amsart} preamble that restates the article's own declarations and macros used by the retained lines, namely the environments \texttt{lem} on separate counters, together with the standard packages those lines need.  The retained environments print in this document as Lemma 1 in order of appearance, whereas the article prints the same items with the numbering of its own full document.  The complete native source of the article as distributed in the arXiv e-print for this exact version is bundled unchanged as \texttt{sources/200264/source0.tex}.
\begin{lem}
We have 
\[
	\theta_{\mathrm{id}\oplus g}(x)(i\cdot l)\otimes v
	=\theta_{\mathrm{id}\oplus g}(x)(i)\cdot l\otimes v 
\]
for $x\in \mathbb{L}_{n}$, $i\in \mathcal{I}_{\mathbb{L}_{n+1}}$, $l\in U(\mathbb{L}_{n+1})$, and $v\in V\otimes_{K}K_{\lambda}$.
\end{lem} 

\begin{proof}
	It is enough to show that 
	\[
		\theta_{\mathrm{id}\oplus g}(x)(l)\cdot v=0.
	\]
Since elements of $U(\mathbb{L}_{n+1})$ are linear combinations of 
products 
of elements of $\mathbb{L}_{n+1}$,
it suffices to show the above equality for 
$l\in \mathbb{L}_{n+1}$.
Moreover since $\theta_{\mathrm{id}\oplus g}|_{\mathbb{L}_{n}}$ coincides with the natural inclusion 
$
\mathbb{L}_{n}\hookrightarrow\mathfrak{P}_{n+1},
$
the equation 
\begin{equation}\label{braidvanish}
	\theta(A_{i,n+1})(l)\cdot v=0
\end{equation}
deduces the above required equation. 

We may assume that 
$l$ is of the form 
$l=[x_{i_{k}},[x_{i_{k-1}},[\ldots,[x_{i_2},x_{i_1}]\ldots]]]$
for $k\ge 1$ where we set $l=x_{i_1}$ if $k=1$,
and show the equation \eqref{braidvanish}
by induction on $k$. 

If $k=1$, we have
\[
\theta(A_{i,n+1})(l)=\theta(A_{i,n+1})(x_{i_1})=
\begin{cases}
\pm [x_{i},x_{n+1}] & \text{if $i_{1}=i$ or $i_{1}=n+1$},\\
0 & \text{otherwise}.
\end{cases}
\]
Since $x_{n+1}\cdot v=\lambda v$, 
we have $[x_{i},x_{n+1}]\cdot v=x_{i}\cdot (\lambda v)-\lambda (x_{i}\cdot v)=0$,
which shows the equation \eqref{braidvanish} in the case $k=1$.

Suppose that the equation \eqref{braidvanish}
holds for $l_{k-1}=[x_{i_{k-1}},[\ldots,[x_{i_2},x_{i_1}]\ldots]]$,
namely, $\theta(A_{i,n+1})(l_{k-1})\cdot v=0$
for any $v\in V\otimes_{K}K_{\lambda}$.
Then we have
\begin{align*}
\theta(A_{i,n+1})(l)\cdot v&=\theta(A_{i,n+1})([x_{i_{k}},l_{k-1}])\cdot v\\
&=
[\theta(A_{i,n+1})(x_{i_{k}}),l_{k-1}]\cdot v
+ [x_{i_{k}},\theta(A_{i,n+1})(l_{k-1})]\cdot v.
\end{align*}
By the induction hypothesis,
$[x_{i_{k}},\theta(A_{i,n+1})(l_{k-1})]\cdot v=0$.
Also the equation 
\[
\theta(A_{i,n+1})(x_{i_k})=
\begin{cases}
\pm [x_{i},x_{n+1}] & \text{if $i_{k}=i$ or $i_{k}=n+1$},\\
0 & \text{otherwise}
\end{cases}
\]
implies 
the first term $[\theta(A_{i,n+1})(x_{i_{k}}),l_{k-1}]\cdot v$ is also zero, 
since $[x_{i},x_{n+1}]\cdot v=0$ for all $v\in V\otimes_{K}K_{\lambda}$.
Thus we are done.
\end{proof}

\end{document}
